\section{Экономическое обоснование}

Актуальность автоматизации процесса составления расписания подтверждается не только внутренними затратами учебного заведения, но и запросом студенческого сообщества. Во «ВКонтакте» размещена петиция «Расписание за неделю» \cite{vk:schedule_petition}, в которой обучающиеся МГТУ «СТАНКИН» фиксируют потребность в своевременной и полной публикации расписания занятий на неделю вперёд, а не фрагментарном или запоздалом доведении информации до студентов и преподавателей. Подобные обращения свидетельствуют о том, что качество и оперативность расписания напрямую влияют на удовлетворённость участников образовательного процесса и воспринимаются как значимый фактор организации учебной деятельности вуза.

С экономической точки зрения такой запрос означает необходимость сократить цикл подготовки и пересборки расписания, снизить число ошибок при ручном согласовании и обеспечить единый канал публикации результата. Разрабатываемое программное средство направлено на устранение этих проблем и тем самым отвечает как организационным потребностям УМУ, так и ожиданиям контингента, выраженным в публичной инициативе.

\subsection{Анализ аналогов и обоснование собственного решения}

Для оценки целесообразности разработки и внедрения программного средства поддержки генерации расписания учебных занятий необходимо рассмотреть существующие подходы к решению данной задачи в образовательных организациях.

В МГТУ «СТАНКИН» на момент выполнения выпускной квалификационной работы расписание на семестр формируется вручную сотрудниками УМУ \cite{stankin:schedule}. Данный процесс является базовым аналогом: он не требует лицензий на программное обеспечение, но связан с высокими трудозатратами, риском ошибок при согласовании ограничений и длительным циклом пересборки расписания при изменении исходных данных.

На российском рынке представлены специализированные системы автоматизированного составления расписания для вузов.

\textbf{Галактика РУЗ} \cite{galaktika:ruz} --- комплексное решение для управления учебным процессом в высших учебных заведениях. Система поддерживает автоматическое формирование опорного расписания, контроль параметризуемых ограничений, публикацию результатов на сайте вуза и в мобильных приложениях. Внедрение предполагает закупку лицензий и сопутствующей инфраструктуры.

\textbf{1С:Автоматизированное составление расписания. Университет} \cite{1c:schedule} --- продукт экосистемы 1С для автоматизированного, ручного и смешанного составления расписаний с учётом нагрузки преподавателей, аудиторного фонда и учебных планов. Типовое внедрение включает лицензирование, настройку под регламенты вуза и интеграцию с другими модулями 1С.

\textbf{ТАНДЕМ.Университет} (модуль «Расписание») \cite{tandem:schedule} --- модуль корпоративной системы управления вузом. Обеспечивает автоматическое и ручное составление расписания, фиксацию ограничений преподавателей и требований к аудиториям, поддержку индивидуальных образовательных траекторий.

\textbf{Экспресс-расписание ВУЗ} \cite{express:schedule} --- программа для автоматического составления расписания в локальной или сетевой конфигурации. Поддерживает работу с подгруппами, планирование отсутствия преподавателей и публикацию расписания на сайте образовательной организации.

Зарубежные системы \textbf{UniTime} и \textbf{TimeEdit}, широко применяемые в зарубежных вузах, рассматриваются в обзорной литературе по задаче составления расписаний \cite{vrielink:patat2016}, однако их внедрение в российский вуз сопряжено с адаптацией нормативной базы, локализацией и интеграцией с отечественными информационными системами.

Разрабатываемое в рамках данной выпускной квалификационной работы программное средство не претендует на полную замену корпоративных ERP-систем, но ориентировано на узкую задачу --- генерацию и оптимизацию расписания с учётом правил МГТУ «СТАНКИН», предпочтений преподавателей, «компоновки вуза» и публикации результата через веб-интерфейс и PWA-клиент. Использование открытого программного стека (Python, FastAPI, Polars, SQLite, Nuxt) снижает стоимость владения по сравнению с коммерческими аналогами.

На основе описания функциональности перечисленных решений проведён сравнительный анализ (см. \tabref{tab:analogs}).

\begin{table}[H]
	\centering
	\caption{Сравнительная таблица функций аналогов и целевого программного продукта}
	\label{tab:analogs}
	\begin{tabularx}{\textwidth}{|>{\raggedright\arraybackslash}X|>{\centering\arraybackslash}p{2.3cm}|>{\centering\arraybackslash}p{2.5cm}|>{\centering\arraybackslash}p{2.5cm}|}
		\hline
		\textbf{Критерий} &
		\textbf{Ручной процесс} &
		\textbf{Коммерч.\ системы} &
		\textbf{Продукт ВКР} \\
		\hline
		Автогенерация с учётом нормативов & $-$ & $+$ & $+$ \\
		\hline
		Учёт предпочтений преподавателей & $\pm$ & $+$ & $+$ \\
		\hline
		Оптимизация перемещений (компоновка) & $-$ & $\pm$ & $+$ \\
		\hline
		Адаптация под правила вуза & $+$ & $\pm$ & $+$ \\
		\hline
		Низкая стоимость владения & $+$ & $-$ & $+$ \\
		\hline
		Веб-доступ и PWA & $\pm$ & $+$ & $+$ \\
		\hline
	\end{tabularx}
\end{table}

Целевой программный продукт превосходит ручной процесс по скорости пересборки расписания и снижению числа конфликтов ограничений. По сравнению с коммерческими ERP-модулями он обеспечивает специализацию под регламенты МГТУ «СТАНКИН» и контролируемую совокупную стоимость владения за счёт открытого программного обеспечения. Таким образом, разработка собственного решения экономически оправдана как альтернатива как полностью ручному составлению, так и закупке дорогостоящего типового продукта без гарантии покрытия всех специфических требований вуза.

\subsection{Оценка затрат проекта}

Для оценки затрат на создание и внедрение программного средства определены следующие ключевые ресурсы:

\begin{itemize}
	\item разработчик (full-stack, автор ВКР) --- проектирование, реализация серверной и клиентской частей, алгоритмов генерации;
	\item специалист по внедрению и сопровождению (системный администратор) --- развёртывание в контуре вуза, обучение пользователей, техническая поддержка;
	\item программное обеспечение --- операционная система, среда разработки, серверные компоненты;
	\item инфраструктура и оборудование --- сервер приложений и рабочее место администратора.
\end{itemize}

Затраты на разработку оцениваются ретроспективно по этапам, отражённым в главах~3--5 выпускной квалификационной работы, с учётом доработок после защиты. Затраты на внедрение и эксплуатацию относятся к первому году использования системы в МГТУ «СТАНКИН».

Для расчёта приняты следующие ставки оплаты труда: разработчик --- 1000~руб./ч, системный администратор --- 500~руб./ч

Расчёт прямых трудозатрат приведён в \tabref{tab:labor}.

\begin{table}[H]
	\centering
	\caption{Расчёт прямых трудозатрат}
	\label{tab:labor}
	\begin{tabularx}{\textwidth}{|>{\raggedright\arraybackslash}X|r|>{\raggedright\arraybackslash}p{2.6cm}|r|}
		\hline
		\textbf{Наименование работ} &
		\textbf{Часы} &
		\textbf{Специалист} &
		\textbf{Затраты, руб.} \\
		\hline
		Проектирование архитектуры и БД (гл.~3) & 80 & Разработчик & 80\,000 \\
		\hline
		Выбор и настройка стека (гл.~4) & 40 & Разработчик & 40\,000 \\
		\hline
		Парсер и наполнение БД (гл.~5) & 60 & Разработчик & 60\,000 \\
		\hline
		Сервер: WFC, оптимизация, API (гл.~5) & 140 & Разработчик & 140\,000 \\
		\hline
		Клиент Nuxt, PWA (гл.~5) & 80 & Разработчик & 80\,000 \\
		\hline
		Тестирование и отладка (гл.~5) & 60 & Разработчик & 60\,000 \\
		\hline
		Доработки после защиты ВКР & 40 & Разработчик & 40\,000 \\
		\hline
		\textbf{Итого разработка} & \textbf{500} & & \textbf{500\,000} \\
		\hline
		Развёртывание в контуре вуза & 60 & Администратор & 30\,000 \\
		\hline
		Обучение сотрудников УМУ & 40 & Администратор & 20\,000 \\
		\hline
		Поддержка ($10$~ч/мес $\times 12$) & 120 & Администратор & 60\,000 \\
		\hline
		\textbf{Итого внедрение и эксплуатация} & \textbf{220} & & \textbf{110\,000} \\
		\hline
	\end{tabularx}
\end{table}

Общая сумма расходов на трудозатраты составляет 610\,000~руб., общая трудоёмкость --- 720~чел.-часов.

Основная часть программного обеспечения проекта распространяется бесплатно (см. \tabref{tab:software}), что соответствует выбранному в главе~4 стеку технологий.

\begin{table}[H]
	\centering
	\caption{Расчёт затрат на программное обеспечение}
	\label{tab:software}
	\begin{tabularx}{\textwidth}{|>{\raggedright\arraybackslash}X|r|>{\raggedright\arraybackslash}X|}
		\hline
		\textbf{Наименование} & \textbf{Стоимость, руб.} & \textbf{Примечание} \\
		\hline
		Python, FastAPI, Polars, SQLite & 0 & Открытый исходный код \\
		\hline
		Nuxt, Vue, TypeScript, pnpm & 0 & Открытый исходный код \\
		\hline
		Ubuntu Server, Git, VS Code & 0 & Бесплатное ПО \\
		\hline
		Лицензия на антивирус & 5\,000 & Стоимость на 1 год \\
		\hline
	\end{tabularx}
\end{table}

Для размещения серверной части и работы администратора необходимо серверное и клиентское оборудование (см. \tabref{tab:infra}). В отличие от типовых смет внедрения учётных систем, в проекте не закладывается лицензия Windows Server --- используется Linux.

\begin{table}[H]
	\centering
	\caption{Расчёт затрат на инфраструктуру}
	\label{tab:infra}
	\begin{tabularx}{\textwidth}{|>{\raggedright\arraybackslash}p{2.5cm}|>{\raggedright\arraybackslash}X|c|r|}
		\hline
		\textbf{Часть} & \textbf{Наименование} & \textbf{Кол-во, ед.} & \textbf{Стоимость, руб.} \\
		\hline
		Серверная & Сервер приложений (8 ядер, 32~ГБ ОЗУ, SSD) & 1 & 100\,000 \\
		\hline
		Клиентская & ПК администратора (16~ГБ ОЗУ) & 1 & 65\,000 \\
		\hline
		Серверная & Сетевое оборудование & 1 & 2\,000 \\
		\hline
		\multicolumn{3}{|r|}{\textbf{Итого}} & \textbf{167\,000} \\
		\hline
	\end{tabularx}
\end{table}

Амортизация оборудования рассчитывается линейным методом по \formref{eq:amort-year}. Срок полезного использования принят равным 7~годам:
\begin{equation}
	A_{\mathrm{год}} = \frac{C_{\mathrm{обор}}}{T_{\mathrm{сл}}},
	\label{eq:amort-year}
\end{equation}
\formwhere{где $A_{\mathrm{год}}$~--- годовая сумма амортизации, руб.;
	\formwherenext{$C_{\mathrm{обор}}$}{первоначальная стоимость оборудования, руб.};
	\formwherenext{$T_{\mathrm{сл}}$}{срок полезного использования, лет}.}

Годовая сумма амортизации составляет $167\,000 / 7 \approx 24\,000$~руб. Для периода внедрения в первый год (5~месяцев) в смету включена доля амортизации 15\,000~руб. (см. \tabref{tab:depreciation}).

\begin{table}[H]
	\centering
	\caption{Расчёт затрат на амортизацию (годовые показатели)}
	\label{tab:depreciation}
	\begin{tabularx}{\textwidth}{|>{\raggedright\arraybackslash}X|r|r|}
		\hline
		\textbf{Наименование} & \textbf{Стоимость, руб.} & \textbf{Амортизация в год, руб.} \\
		\hline
		Сервер приложений & 100\,000 & 14\,286 \\
		\hline
		ПК администратора & 65\,000 & 9\,286 \\
		\hline
		Сетевое оборудование & 2\,000 & 286 \\
		\hline
		\textbf{Итого} & \textbf{167\,000} & \textbf{24\,000} \\
		\hline
	\end{tabularx}
\end{table}

Затраты на электроэнергию за год оцениваются по \formref{eq:energy-cost}:
\begin{equation}
	Z_{\mathrm{эл}} = P \cdot t \cdot c,
	\label{eq:energy-cost}
\end{equation}
\formwhere{где $Z_{\mathrm{эл}}$~--- затраты на электроэнергию за год, руб.;
	\formwherenext{$P$}{потребляемая мощность, кВт};
	\formwherenext{$t$}{время работы оборудования в год, ч};
	\formwherenext{$c$}{тариф, руб./(кВт$\cdot$ч)}.}

Для сервера и рабочего места администратора суммарная мощность принята равной 0{,}5~кВт при 1600~ч работы в год и тарифе 5{,}47~руб./(кВт$\cdot$ч):
$Z_{\mathrm{эл}} = 0{,}5 \cdot 1600 \cdot 5{,}47 \approx 4\,400$~руб./год.

Затраты на ремонт и обслуживание оборудования оцениваются в 10\,\% от его стоимости в год: $167\,000 \cdot 0{,}1 = 16\,700$~руб.

Исходя из проведённых вычислений, сформирована общая смета затрат на первый год внедрения (см. \tabref{tab:budget}).

\begin{table}[H]
	\centering
	\caption{Общая смета затрат на проект (первый год)}
	\label{tab:budget}
	\begin{tabularx}{\textwidth}{|>{\raggedright\arraybackslash}X|r|}
		\hline
		\textbf{Статья затрат} & \textbf{Сумма, руб.} \\
		\hline
		Трудозатраты разработки (ВКР и доработки) & 500\,000 \\
		\hline
		Трудозатраты внедрения и поддержки & 110\,000 \\
		\hline
		Программное обеспечение & 5\,000 \\
		\hline
		Инфраструктура & 167\,000 \\
		\hline
		Амортизация (доля за период внедрения) & 15\,000 \\
		\hline
		Электроэнергия и ремонт & 21\,100 \\
		\hline
		\textbf{Итого} & \textbf{818\,100} \\
		\hline
	\end{tabularx}
\end{table}

\subsection{Обоснование экономической выгоды проекта}

Внедрение программного средства поддержки генерации расписания обеспечит следующие экономические преимущества для МГТУ «СТАНКИН»:

\begin{itemize}
	\item сокращение ручного труда сотрудников УМУ при пересборке расписания на семестр;
	\item снижение числа конфликтов по аудиториям, нагрузке преподавателей и «окнам» в расписании;
	\item ускорение публикации актуального расписания через веб-интерфейс и PWA вместо разрозненных файлов;
	\item учёт предпочтений преподавателей и оптимизация перемещений между корпусами без многократных ручных согласований.
\end{itemize}

Для количественной оценки эффекта используются допущения. Базовые трудозатраты УМУ на ручное составление расписания оцениваются так: 3~сотрудника $\times$ 120~ч $\times$ 500~руб./ч $=$ 180\,000~руб. за семестр, или 360\,000~руб. в год (два семестра).

Автоматизация позволяет сократить трудозатраты на подготовку и корректировку расписания на 35--40\,\%. При коэффициенте сокращения $k = 0{,}375$ годовая экономия определяется \formref{eq:year-savings}:
\begin{equation}
	E_{\mathrm{год}} = Z_{\mathrm{баз}} \cdot k = 360\,000 \cdot 0{,}375 = 135\,000~\text{руб.}
	\label{eq:year-savings}
\end{equation}
\formwhere{где $E_{\mathrm{год}}$~--- годовая экономия трудозатрат, руб.;
	\formwherenext{$Z_{\mathrm{баз}}$}{базовые годовые трудозатраты УМУ без автоматизации, руб.};
	\formwherenext{$k$}{коэффициент сокращения трудозатрат}.}

Сравнение с закупкой коммерческого решения для вуза сопоставимого масштаба (по аналогии с оценками внедрения типовых учётных систем) показывает, что совокупные затраты на лицензию, кастомизацию и внедрение могут составить 1{,}5--2{,}5~млн~руб. Разработка и внедрение собственного средства по смете \tabref{tab:budget} обходятся примерно в 818\,100~руб., то есть прямая экономия на капитальных затратах (CAPEX) по сравнению с коммерческим аналогом составляет порядка 0{,}7--1{,}7~млн~руб. при сопоставимом базовом функционале генерации расписания.

Срок окупаемости проекта по операционной экономии УМУ определяется \formref{eq:payback}:
\begin{equation}
	T_{\mathrm{ок}} = \frac{Z_{\mathrm{проект}}}{E_{\mathrm{год}}} = \frac{818\,100}{135\,000} \approx 6{,}1~\text{года}.
	\label{eq:payback}
\end{equation}
\formwhere{где $T_{\mathrm{ок}}$~--- срок окупаемости проекта, лет;
	\formwherenext{$Z_{\mathrm{проект}}$}{затраты на проект по смете, руб.};
	\formwherenext{$E_{\mathrm{год}}$}{годовая экономия трудозатрат, руб.}.}

При учёте отказа от закупки коммерческой лицензии совокупный экономический эффект превышает затраты на проект уже в первый год эксплуатации.

Дополнительно повышается качество расписания: уменьшение «окон», равномерность нагрузки и учёт «компоновки вуза» снижают косвенные потери времени студентов и преподавателей, не полностью отражённые в расчёте $E_{\mathrm{год}}$.

\subsection{Вывод по экономическому обоснованию}
Разработка и внедрение программного средства поддержки генерации расписания для МГТУ «СТАНКИН» экономически целесообразны как альтернатива полностью ручному процессу и закупке дорогостоящих типовых ERP-модулей. Низкая совокупная стоимость владения обеспечивается использованием открытого программного обеспечения; операционная экономия формируется за счёт сокращения трудозатрат УМУ. Ограничением остаётся необходимость доработки интеграции с корпоративными системами вуза (ЭОС, учебные планы), отражённая в главе~5. Для небольших подразделений с минимальным объёмом расписания внедрение может быть менее выгодным; для вуза уровня МГТУ «СТАНКИН» проект является обоснованным с точки зрения соотношения затрат и эффекта.
